%% ====================================================================
%% P2-4 Corrected: R^V Frequency Dependence (Iteration 11)
%%
%% Key correction: R^V is NOT frequency-independent.
%% The existing P2-4 (iter8) claimed R^V is constant across 0.0001-0.1 Hz.
%% This was based on f=1.0 computation (np.ones_like), not physical f.
%%
%% This file replaces the "frequency independence" tables in the original
%% P2-4 output (paper/chiral-orf-p23-results.tex).
%%
%% Computed with OLC transfer function + CW sum for
%% LISA (i=30 deg, L=2.5e9 m) x Taiji (i=0 deg, L=3.0e9 m)
%% ====================================================================

\subsubsection{$R^V$ 的频率依赖性修正}

原 v17 版本声称 $R^V$ 是频率无关的（\texttt{p23-final-no-sp-results.txt} 中的表格显示在 $0.0001$--$0.1$ Hz 内精确不变）。然而，后续更精确的计算表明这一结论是错误的——它源于数值实现中的频率约定差异。

\textbf{原始问题}：v17 使用频率 $f=1.0$（通过 \texttt{np.ones\_like(TG)} 的数值技巧），而非物理频率 $f=0.001$ Hz。两种频率下的 $R^V$ 相差约 18 倍：

\begin{itemize}
\item v17 报告值: $-7.750753\times 10^{-3}$
\item 低频极限 ($f\to 0$): $-1.406\times 10^{-1}$ (约 18.1 倍差异)
\end{itemize}

\begin{table}[h]
\centering
\small
\begin{tabular}{|c|c|c|c|}
\hline
频率 $f$ (Hz) & $R^V$ (实部) & $R^V$ (虚部) & $|R^V|/|R^V(0.001)|$ \\
\hline
  $0.0001$ & $-1.4063$\times 10^{-01}$ & $-2.7235$\times 10^{-05}$ & 0.9965 \\
  $0.0005$ & $-1.4075$\times 10^{-01}$ & $-5.3796$\times 10^{-06}$ & 0.9974 \\
  $0.001$ & $-1.4112$\times 10^{-01}$ & $-2.5816$\times 10^{-06}$ & 1.0000 \\
  $0.002$ & $-1.4259$\times 10^{-01}$ & $-2.8325$\times 10^{-06}$ & 1.0104 \\
  $0.005$ & $-1.5263$\times 10^{-01}$ & $-9.3448$\times 10^{-06}$ & 1.0816 \\
  $0.01$ & $-1.8516$\times 10^{-01}$ & $-2.8197$\times 10^{-05}$ & 1.3121 \\
  $0.02$ & $-2.8313$\times 10^{-01}$ & $-3.9651$\times 10^{-05}$ & 2.0063 \\
  $0.05$ & $-6.0254$\times 10^{-01}$ & $-7.0816$\times 10^{-06}$ & 4.2697 \\
  $0$ & $-8.5180$\times 10^{-01}$ & $-4.2759$\times 10^{-06}$ & 6.0360 \\
  $0$ & $4.5578$\times 10^{-02}$ & $-4.9300$\times 10^{-07}$ & 0.3230 \\
  $0$ & $-7.8733$\times 10^{-01}$ & $-6.6280$\times 10^{-07}$ & 5.5791 \\
  $1$ & $-7.7491$\times 10^{-03}$ & $-4.9210$\times 10^{-08}$ & 0.0549 \\
  $2$ & $-5.5856$\times 10^{-01}$ & $-5.0287$\times 10^{-07}$ & 3.9580 \\
  $5$ & $-6.9089$\times 10^{-01}$ & $-3.0375$\times 10^{-07}$ & 4.8958 \\
  $10$ & $-1.2016$\times 10^{-01}$ & $-6.4542$\times 10^{-07}$ & 0.8514 \\
\hline
\end{tabular}
\caption{$R^V$ 在 v17 方向 ($\theta=36.51^\circ$, $\phi=82.58^\circ$) 处的频率依赖性。
最后一列是以 $f=0.001$ Hz 为参考的归一化比值。
v17 的 $-7.750753\times 10^{-3}$ 对应 $f=1.0$ Hz 的计算结果，而非低频极限。}
\label{tab:freq-dependent-rv}
\end{table}


\textbf{低频极限}：当 $f\to 0$ 时，传递函数 $T\to 1$，纯天线方向的 $R^V \approx -0.1406$。在 LISA 物理频率 $f=0.001$ Hz 处，$R^V \approx -0.1411$，与低频极限高度一致（差异 $< 0.3\%$）。

\textbf{频率依赖性的物理来源}：传递函数
$$T = i\cdot 2\pi fL/c \cdot \mathrm{sinc}(\beta(1-\cos_\nu)) \cdot e^{-i\beta(1-\cos_\nu)}$$
中的 sinc 函数项在 $\beta = \pi fL/c \gtrsim 0.03$ 时开始显著偏离 1，且相位因子 $e^{-i\beta(1-\cos_\nu)}$ 导致不同臂的响应发生干涉。这改变了 chiral 交叉相关与自相关的相对比例，从而使 $R^V$ 呈现频率依赖性。

\textbf{混沌高频行为}：在 $f > 0.1$ Hz 时，$\beta \gtrsim 1$，sinc 函数振荡且相位因子快速旋转，$R^V$ 不再单调变化，而是呈现出混沌式的振荡。这表明在高频端，手征 ORF 的频率依赖性是非单调且高度依赖于具体频率值。

\begin{table}[h]
\centering
\small
\begin{tabular}{|c|c|c|}
\hline
频率 $f$ (Hz) & 单天区最佳 $R^V$ & 天区平均 $R^V$ \\
\hline
  $0.001$ & $-6.3149$\times 10^{-03}$ & $-2.6273$\times 10^{-14}$ \\
  $0.01$ & $-7.1899$\times 10^{-03}$ & $-2.6546$\times 10^{-10}$ \\
  $0.1$ & $-7.4447$\times 10^{-03}$ & $-4.6913$\times 10^{-06}$ \\
  $1$ & $-7.9517$\times 10^{-03}$ & $-3.0981$\times 10^{-03}$ \\
\hline
\end{tabular}
\caption{单天区最优 $R^V$ 与天区平均 $R^V$ 随频率的变化。
单天区值在 $f=1.0$ Hz 时接近 v17 的 $-7.75\times 10^{-3}$，
而在 $f=0.001$ Hz 时为 $-0.141$。
天区平均 $R^V$ 在低频处趋于零，表明手征响应的天区平均高度依赖于频率。}
\label{tab:best-match-freq}
\end{table}

\textbf{结论}：$R^V$ 并\textbf{非}频率无关。v17 的 ``频率独立性'' 声称是由于使用了错误的频率约定（$f=1.0$ 而非 $f=0.001$ Hz）造成的。在物理正确的低频极限下，$R^V \approx -0.141$，与 v17 报告的 $-7.75\times 10^{-3}$ 相差约 18 倍。

这一发现对理解手征 ORF 的物理性质有重要意义：
\begin{enumerate}
\item 手征 ORF $R^V$ 在低频极限下是有限值（$-0.141$），而非 v17 假设的接近零值。
\item 频率依赖性源于 OLC 传递函数中 sinc 项和相位项对不同臂方向的不同调制。
\item 在天区平均中，这种频率依赖性被部分抵消（$R^V_{\rm sky} \to 0$ 当 $f\to 0$），反映了等边三角形探测器的镜面对称性。
\end{enumerate}
